\documentclass[10pt]{extarticle}
\usepackage[a4paper,margin=1.5cm]{geometry}
\usepackage[T1]{fontenc}
\usepackage{amsmath,amssymb}
\usepackage{graphicx}
\usepackage[export]{adjustbox}
\usepackage{multicol}
\usepackage{enumitem}
\usepackage{booktabs,array,tabularx}
\setlength{\columnsep}{0.7cm}
\setlength{\parindent}{0pt}
\setlength{\parskip}{3pt}
\pagestyle{empty}
\makeatletter
\newcommand{\qfigdef}[4]{\@namedef{qf@file@#1}{#2}\@namedef{qf@w@#1}{#3}\@namedef{qf@wide@#1}{#4}}
\newcommand{\qf@check}[1]{\@ifundefined{qf@file@#1}{\PackageError{score2}{Unknown figure id #1}{}}{}}
\newcommand{\qf@img}[4][]{\includegraphics[width=#2\dimexpr\@nameuse{qf@w@#4}pt\relax,max width=#3,max height=0.45\textheight,#1]{\@nameuse{qf@file@#4}}}
\newif\ifqf@full
\newif\ifqf@boxed
\newsavebox{\qf@box}
\newcommand{\qf@open}{\global\qf@fulltrue\par\global\qf@boxedtrue
  \setbox\qf@box\vbox\bgroup\hsize\textwidth\linewidth\textwidth}
\newcommand{\qf@emit}[1]{\par\egroup\global\qf@boxedfalse
  \ifdim\dimexpr\ht\qf@box+\dp\qf@box\relax<#1\textheight\noindent\usebox{\qf@box}\par\else\unvbox\qf@box\fi}
\newcommand{\qfullsplit}{\ifqf@boxed\qf@emit{0.85}\fi}
\newcommand{\qf@close}{\ifqf@boxed\qf@emit{0.45}\fi\global\qf@fullfalse}
\newcommand{\qfullbegin}{\end{multicols}\qf@open}
\newcommand{\qfullend}{\qf@close\begin{multicols}{2}}
\newcommand{\qfullstart}{\qf@open}
\newcommand{\qfulldone}{\qf@close}
\newcommand{\qfig}[2][1]{\qf@check{#2}%
  \ifqf@full\qf@blk{#1}{#2}\else
  \ifnum\@nameuse{qf@wide@#2}=1 \expandafter\@firstoftwo\else\expandafter\@secondoftwo\fi
  {\end{multicols}\begin{center}\qf@img{#1}{\textwidth}{#2}\end{center}\begin{multicols}{2}}%
  {\qf@blk{#1}{#2}}\fi}
\newcommand{\qf@blk}[2]{\par\smallskip{\centering\qf@img{#1}{\linewidth}{#2}\par}\smallskip}
\newcommand{\qfigi}[2][1]{\qf@check{#2}\qf@img[valign=c]{#1}{\linewidth}{#2}}
\makeatother
\newcommand{\papersec}[1]{\par\bigskip{\centering\Large\bfseries PAPER-#1\par}\medskip}
\newcommand{\partsec}[2]{\par\medskip\noindent{\large\bfseries PART-#1: #2}\par\smallskip}
\newcommand{\sect}[1]{\par\smallskip\noindent\textbf{#1}\par}
\newcommand{\qn}[1]{\par\smallskip\noindent\textbf{#1.}~}
\newcommand{\paraq}[2]{\par\smallskip\noindent\textbf{Paragraph for Questions #1 and #2}\par}
\newenvironment{qoptions}{\begin{itemize}[leftmargin=1.8em,itemsep=1pt,topsep=2pt,label={}]}{\end{itemize}}
\newcommand{\opt}[2]{\item[(#1)] #2}
% figure registry (generated by score2_tools.py finalize)
\qfigdef{p1_phy_s1_q01_1}{Score-I - Test 04 (14 Jan 2022)-img/p1_phy_s1_q01_1.png}{131.8}{0}
\qfigdef{p1_phy_s1_q02_A}{Score-I - Test 04 (14 Jan 2022)-img/p1_phy_s1_q02_A.png}{134.5}{0}
\qfigdef{p1_phy_s1_q03_1}{Score-I - Test 04 (14 Jan 2022)-img/p1_phy_s1_q03_1.png}{84.0}{0}
\qfigdef{p1_phy_s2_para7_8_1}{Score-I - Test 04 (14 Jan 2022)-img/p1_phy_s2_para7_8_1.png}{237.5}{0}
\qfigdef{p1_phy_s2_q07_A}{Score-I - Test 04 (14 Jan 2022)-img/p1_phy_s2_q07_A.png}{106.3}{0}
\qfigdef{p1_phy_s2_q07_B}{Score-I - Test 04 (14 Jan 2022)-img/p1_phy_s2_q07_B.png}{106.8}{0}
\qfigdef{p1_phy_s2_q07_C}{Score-I - Test 04 (14 Jan 2022)-img/p1_phy_s2_q07_C.png}{106.3}{0}
\qfigdef{p1_phy_s2_q07_D}{Score-I - Test 04 (14 Jan 2022)-img/p1_phy_s2_q07_D.png}{107.2}{0}
\qfigdef{p1_phy_s2_q08_A}{Score-I - Test 04 (14 Jan 2022)-img/p1_phy_s2_q08_A.png}{106.0}{0}
\qfigdef{p1_phy_s2_q08_B}{Score-I - Test 04 (14 Jan 2022)-img/p1_phy_s2_q08_B.png}{105.1}{0}
\qfigdef{p1_phy_s2_q08_C}{Score-I - Test 04 (14 Jan 2022)-img/p1_phy_s2_q08_C.png}{106.3}{0}
\qfigdef{p1_phy_s2_q08_D}{Score-I - Test 04 (14 Jan 2022)-img/p1_phy_s2_q08_D.png}{105.6}{0}
\qfigdef{p1_phy_s2_para9_10_1}{Score-I - Test 04 (14 Jan 2022)-img/p1_phy_s2_para9_10_1.png}{208.0}{0}
\qfigdef{p1_phy_s2_q09_1}{Score-I - Test 04 (14 Jan 2022)-img/p1_phy_s2_q09_1.png}{96.3}{0}
\qfigdef{p1_phy_s2_q09_2}{Score-I - Test 04 (14 Jan 2022)-img/p1_phy_s2_q09_2.png}{127.2}{0}
\qfigdef{p1_phy_s2_q09_3}{Score-I - Test 04 (14 Jan 2022)-img/p1_phy_s2_q09_3.png}{176.8}{0}
\qfigdef{p1_phy_s2_q09_4}{Score-I - Test 04 (14 Jan 2022)-img/p1_phy_s2_q09_4.png}{84.4}{0}
\qfigdef{p1_phy_s3_q01_1}{Score-I - Test 04 (14 Jan 2022)-img/p1_phy_s3_q01_1.png}{128.8}{0}
\qfigdef{p1_phy_s3_q02_1}{Score-I - Test 04 (14 Jan 2022)-img/p1_phy_s3_q02_1.png}{124.6}{0}
\qfigdef{p1_phy_s3_q05_1}{Score-I - Test 04 (14 Jan 2022)-img/p1_phy_s3_q05_1.png}{122.4}{0}
\qfigdef{p1_phy_s3_q06_1}{Score-I - Test 04 (14 Jan 2022)-img/p1_phy_s3_q06_1.png}{152.3}{0}
\qfigdef{p1_phy_s3_q07_1}{Score-I - Test 04 (14 Jan 2022)-img/p1_phy_s3_q07_1.png}{156.7}{0}
\qfigdef{p1_chem_s1_q02_1}{Score-I - Test 04 (14 Jan 2022)-img/p1_chem_s1_q02_1.png}{60.4}{0}
\qfigdef{p1_chem_s1_q02_A}{Score-I - Test 04 (14 Jan 2022)-img/p1_chem_s1_q02_A.png}{65.8}{0}
\qfigdef{p1_chem_s1_q02_B}{Score-I - Test 04 (14 Jan 2022)-img/p1_chem_s1_q02_B.png}{72.2}{0}
\qfigdef{p1_chem_s1_q02_C}{Score-I - Test 04 (14 Jan 2022)-img/p1_chem_s1_q02_C.png}{64.2}{0}
\qfigdef{p1_chem_s1_q02_E}{Score-I - Test 04 (14 Jan 2022)-img/p1_chem_s1_q02_E.png}{38.7}{0}
\qfigdef{p1_chem_s1_q02_D}{Score-I - Test 04 (14 Jan 2022)-img/p1_chem_s1_q02_D.png}{86.3}{0}
\qfigdef{p1_chem_s2_para9_10_1}{Score-I - Test 04 (14 Jan 2022)-img/p1_chem_s2_para9_10_1.png}{245.2}{0}
\qfigdef{p1_chem_s3_q05_1}{Score-I - Test 04 (14 Jan 2022)-img/p1_chem_s3_q05_1.png}{209.1}{0}
\qfigdef{p1_chem_s3_q06_1}{Score-I - Test 04 (14 Jan 2022)-img/p1_chem_s3_q06_1.png}{97.2}{0}
\qfigdef{p2_phy_s1_q01_1}{Score-I - Test 04 (14 Jan 2022)-img/p2_phy_s1_q01_1.png}{89.1}{0}
\qfigdef{p2_phy_s1_q02_1}{Score-I - Test 04 (14 Jan 2022)-img/p2_phy_s1_q02_1.png}{132.8}{0}
\qfigdef{p2_phy_s1_q03_1}{Score-I - Test 04 (14 Jan 2022)-img/p2_phy_s1_q03_1.png}{104.4}{0}
\qfigdef{p2_phy_s1_q04_1}{Score-I - Test 04 (14 Jan 2022)-img/p2_phy_s1_q04_1.png}{119.6}{0}
\qfigdef{p2_phy_s1_q05_1}{Score-I - Test 04 (14 Jan 2022)-img/p2_phy_s1_q05_1.png}{148.0}{0}
\qfigdef{p2_phy_s1_q06_1}{Score-I - Test 04 (14 Jan 2022)-img/p2_phy_s1_q06_1.png}{145.4}{0}
\qfigdef{p2_phy_s1_q07_1}{Score-I - Test 04 (14 Jan 2022)-img/p2_phy_s1_q07_1.png}{154.3}{0}
\qfigdef{p2_phy_s1_q09_1}{Score-I - Test 04 (14 Jan 2022)-img/p2_phy_s1_q09_1.png}{169.4}{0}
\qfigdef{p2_phy_s1_q10_1}{Score-I - Test 04 (14 Jan 2022)-img/p2_phy_s1_q10_1.png}{152.1}{0}
\qfigdef{p2_phy_s1_q12_1}{Score-I - Test 04 (14 Jan 2022)-img/p2_phy_s1_q12_1.png}{150.1}{0}
\qfigdef{p2_phy_s1_q16_A}{Score-I - Test 04 (14 Jan 2022)-img/p2_phy_s1_q16_A.png}{61.7}{0}
\qfigdef{p2_phy_s1_q16_B}{Score-I - Test 04 (14 Jan 2022)-img/p2_phy_s1_q16_B.png}{61.7}{0}
\qfigdef{p2_phy_s1_q16_C}{Score-I - Test 04 (14 Jan 2022)-img/p2_phy_s1_q16_C.png}{61.7}{0}
\qfigdef{p2_phy_s1_q16_D}{Score-I - Test 04 (14 Jan 2022)-img/p2_phy_s1_q16_D.png}{61.7}{0}
\qfigdef{p2_phy_s1_q18_1}{Score-I - Test 04 (14 Jan 2022)-img/p2_phy_s1_q18_1.png}{147.5}{0}
\qfigdef{p2_phy_s1_q19_1}{Score-I - Test 04 (14 Jan 2022)-img/p2_phy_s1_q19_1.png}{76.9}{0}
\qfigdef{p2_phy_s1_q20_A}{Score-I - Test 04 (14 Jan 2022)-img/p2_phy_s1_q20_A.png}{77.9}{0}
\qfigdef{p2_phy_s2_q01_1}{Score-I - Test 04 (14 Jan 2022)-img/p2_phy_s2_q01_1.png}{134.7}{0}
\qfigdef{p2_phy_s2_q02_1}{Score-I - Test 04 (14 Jan 2022)-img/p2_phy_s2_q02_1.png}{145.4}{0}
\qfigdef{p2_phy_s2_q03_1}{Score-I - Test 04 (14 Jan 2022)-img/p2_phy_s2_q03_1.png}{78.2}{0}
\qfigdef{p2_phy_s2_q04_1}{Score-I - Test 04 (14 Jan 2022)-img/p2_phy_s2_q04_1.png}{200.2}{0}
\qfigdef{p2_phy_s2_q09_1}{Score-I - Test 04 (14 Jan 2022)-img/p2_phy_s2_q09_1.png}{154.2}{0}
\qfigdef{p2_phy_s2_q10_1}{Score-I - Test 04 (14 Jan 2022)-img/p2_phy_s2_q10_1.png}{154.5}{0}
\qfigdef{p2_chem_s1_q12_A}{Score-I - Test 04 (14 Jan 2022)-img/p2_chem_s1_q12_A.png}{105.1}{0}
\qfigdef{p2_chem_s1_q12_B}{Score-I - Test 04 (14 Jan 2022)-img/p2_chem_s1_q12_B.png}{105.5}{0}
\qfigdef{p2_chem_s1_q12_C}{Score-I - Test 04 (14 Jan 2022)-img/p2_chem_s1_q12_C.png}{106.1}{0}
\qfigdef{p2_chem_s1_q12_D}{Score-I - Test 04 (14 Jan 2022)-img/p2_chem_s1_q12_D.png}{106.5}{0}
\qfigdef{p2_chem_s1_q15_1}{Score-I - Test 04 (14 Jan 2022)-img/p2_chem_s1_q15_1.png}{29.2}{0}
\qfigdef{p2_chem_s1_q16_1}{Score-I - Test 04 (14 Jan 2022)-img/p2_chem_s1_q16_1.png}{39.4}{0}
\qfigdef{p2_chem_s1_q20_A1}{Score-I - Test 04 (14 Jan 2022)-img/p2_chem_s1_q20_A1.png}{30.3}{0}
\qfigdef{p2_chem_s1_q20_A2}{Score-I - Test 04 (14 Jan 2022)-img/p2_chem_s1_q20_A2.png}{39.4}{0}
\qfigdef{p2_chem_s1_q20_B1}{Score-I - Test 04 (14 Jan 2022)-img/p2_chem_s1_q20_B1.png}{30.4}{0}
\qfigdef{p2_chem_s1_q20_B2}{Score-I - Test 04 (14 Jan 2022)-img/p2_chem_s1_q20_B2.png}{35.0}{0}
\qfigdef{p2_chem_s1_q20_C1}{Score-I - Test 04 (14 Jan 2022)-img/p2_chem_s1_q20_C1.png}{39.0}{0}
\qfigdef{p2_chem_s1_q20_C2}{Score-I - Test 04 (14 Jan 2022)-img/p2_chem_s1_q20_C2.png}{32.6}{0}
\qfigdef{p2_chem_s1_q20_D1}{Score-I - Test 04 (14 Jan 2022)-img/p2_chem_s1_q20_D1.png}{38.5}{0}
\qfigdef{p2_chem_s2_q01_1}{Score-I - Test 04 (14 Jan 2022)-img/p2_chem_s2_q01_1.png}{151.3}{0}
\qfigdef{p2_chem_s2_q02_a}{Score-I - Test 04 (14 Jan 2022)-img/p2_chem_s2_q02_a.png}{49.0}{0}
\qfigdef{p2_chem_s2_q02_b}{Score-I - Test 04 (14 Jan 2022)-img/p2_chem_s2_q02_b.png}{46.3}{0}
\qfigdef{p2_chem_s2_q02_d}{Score-I - Test 04 (14 Jan 2022)-img/p2_chem_s2_q02_d.png}{52.4}{0}
\qfigdef{p2_chem_s2_q06_1}{Score-I - Test 04 (14 Jan 2022)-img/p2_chem_s2_q06_1.png}{65.2}{0}

\begin{document}
\begin{center}{\Large\bfseries JEE (Main+Advanced) Score-I --- Test 04 (14 Jan 2022)}\end{center}
\begin{multicols}{2}
\papersec{1}
\partsec{1}{PHYSICS}
\sect{Section-I(i)}
\qn{1} Consider a long thin conducting (having large number of free electron) hollow cylindrical tube of radius R, thickness t, made up of material of conductivity $\sigma$. It is rotated about it's axis of symmetry with angular acceleration $\alpha$, starting from rest. Take mass of electron to be $m$ \& charge magnitude $e$.
\qfig{p1_phy_s1_q01_1}
\begin{qoptions}
\opt{A}{Due to different behaviour of free electrons and fixed lattice ions, induced EMF will be observed in circumference of tube.}
\opt{B}{Induced current density developed in tube along it's circumference is, $J=\dfrac{mR\sigma\alpha}{e}$.}
\opt{C}{A constant hall effect like potential difference is seen between inner and outer curved surface of cylinder.}
\opt{D}{Magnetic field seen near center of cylinder is $B=\dfrac{\mu_0 mR\sigma\alpha t}{e}$.}
\end{qoptions}
\qn{2} Consider a spiral loop of inner radius a \& outer radius b, made up of N turns of fine close wires with current I. O is center of spiral. P is a point on axis of spiral such that outer \& inner radius makes angle $\dfrac{\pi}{3}$ \& $\dfrac{\pi}{6}$ on that point.
\qfig{p1_phy_s1_q02_A}
\begin{qoptions}
\opt{A}{$B_P=\dfrac{\mu_0 iN}{(b-a)}\left[\ell n\left(\dfrac{\left(2-\sqrt3\right)}{\sqrt3}\right)-\dfrac12\left(\sqrt3-1\right)\right]$}
\opt{B}{$B_P=\dfrac{\mu_0 iN}{2(b-a)}\left[\ell n\left(\dfrac{\left(2+\sqrt3\right)}{\sqrt3}\right)-\dfrac12\left(\sqrt3-1\right)\right]$}
\opt{C}{$B_O=\dfrac{\mu_0 iN}{(b-a)}\ell n\left(\sqrt{\dfrac ba}\right)$}
\opt{D}{$B_O=\dfrac{\mu_0 iN}{(b-a)}\ell n\left(\dfrac ba\right)$}
\end{qoptions}
\qn{3} PP$'$ and QQ$'$ are two fixed infinite long straight wire both parallel to $z$ axis, with current I in each wire, ABCD is a square loop of side $2\ell$. Sides AD and BC are parallel to $z$ axis \& AB is in $xy$ plane and, CD is in plane parallel to $x$-$y$ plane.
\qfig{p1_phy_s1_q03_1}
PP$'$ and QQ$'$ are cutting $x$-axis at $(-a,0,0)$ and $(a,0,0)$ and $2\ell<2a$. ABCD loop is free to rotate about $z$-axis only which is passing in it's plane from center of ABCD loop. Side AB of loop is making angle $\phi$ with $x$-axis. Choose correct statements.
\begin{qoptions}
\opt{A}{Torque on ABCD loop along $z$ axis in situation when $\phi=0$ is zero.}
\opt{B}{Torque on ABCD loop along $z$ axis in situation when $\phi=\dfrac{\pi}{2}$ is zero.}
\opt{C}{Torque on ABCD loop along $z$ axis in situation when $\phi=\dfrac{\pi}{2}$ is $\dfrac{4\mu_0 I^2\ell^2 a}{\pi\left(a^2+\ell^2\right)}$.}
\opt{D}{Torque on ABCD loop along $z$ axis in situation when $\phi=\dfrac{\pi}{2}$ is maximum.}
\end{qoptions}
\qn{4} An alternating current generator has an internal resistance $R_g$ and an internal inductive reactance of magnitude $X_g$. It is used to supply power to a passive load consisting of a resistance $R_L$ (fixed) and a capacitive reactance of magnitude $X_c$. For maximum power to be delivered from the generator to the load the value of $X_c$ is equal to :
\begin{qoptions}
\opt{A}{$R_g$}
\opt{B}{$X_g$}
\opt{C}{Equal and opposite of reactance of generator.}
\opt{D}{Equal in magnitude of impendence of generator.}
\end{qoptions}
\qn{5} A car is moving on AB road Indore from north to south direction, with speed 72 km/hr. An electric bulb is connected between right to left metal body pannel of car. Resistance of bulb is 10$\Omega$, and you can assume resistance in body of car and connecting metal wire is small. If width of car is 1.5 m, \& horizontal component of earth's magnetic field in south to north $B_H=10$ Gauss with dip angle $45^\circ$, then
\begin{qoptions}
\opt{A}{Potential difference accross bulb due to induced current is 30 mV.}
\opt{B}{Current induced in bulb is 3 mA.}
\opt{C}{EMF induced in body of car along it's width is 30 mV.}
\opt{D}{Current induced in bulb is zero.}
\end{qoptions}
\qn{6} In working of moving coil galvanometer (select the correct alternatives)
\begin{qoptions}
\opt{A}{Eddy current help in achieving equilibrium of coil faster.}
\opt{B}{Cylindrical soft iron core helps in make magnetic field radial \& increasing strength of torque on coil.}
\opt{C}{Parallex error needs to minized while taking measurement.}
\opt{D}{Marking on scale are equally spaced for equal change in current.}
\end{qoptions}
\sect{Section-I(ii)}
\paraq{7}{8}
Transcranial magnetic stimulation (TMS) is a nonivasive form of brain stimulation in which a changing magnetic field is used to cause electric current at a specific area of the brain through electromagnetic induction. An electric pulse generator, or stimulator, is connected to a magnetic coil, which in turn is connected to the scalp. The stimulator generates a changing electric current within the coil wich induces a magnetic field; this field then causes a second inductance of inverted electric charge within the brain itself. Transcranial magnetic stimulation is achieved by quickly discharging current from a large capacitor into a coil to produce pulsed magnetic fields between 2 and 3 Tesla in strength. Directing the magnetic field pulse at a targeted area in the brain causes a localized electrical current which can then either depolarize or hyperpolarize neurons at that site.

Some design of circuit for discharing of capacitor through coil uses diode, which is like a current direction controller. Diode allow current in one direction \& not in other.
\qfig{p1_phy_s2_para7_8_1}
Circuit for this design is as shown. When switch is closed charged capacitor discharges through inductor coil (L) then diode D$_2$ allows inductor to discharge through large resistance R. Diode D$_1$ prevents LC circuit to produce current in oppsite direction in inductor.
\qn{7} Correct graph of current in coil (L) v/s time should look like:
\begin{qoptions}
\opt{A}{\qfigi{p1_phy_s2_q07_A}}
\opt{B}{\qfigi{p1_phy_s2_q07_B}}
\opt{C}{\qfigi{p1_phy_s2_q07_C}}
\opt{D}{\qfigi{p1_phy_s2_q07_D}}
\end{qoptions}
\qn{8} Correct graph of induced electric field due to coil v/s time will be :
\begin{qoptions}
\opt{A}{\qfigi{p1_phy_s2_q08_A}}
\opt{B}{\qfigi{p1_phy_s2_q08_B}}
\opt{C}{\qfigi{p1_phy_s2_q08_C}}
\opt{D}{\qfigi{p1_phy_s2_q08_D}}
\end{qoptions}
\paraq{9}{10}
Arrangement shown below has an electron source producing electrons of various speed emitted along line of $x$-axis. Region-I is sufficiently large in dimensions to work as velocity selector. It contain uniform \& constant with time, magnetic field of strength B and electric field of strength E, which are along $-\hat k$ and $-\hat j$ direction as shown.

Region-II contain uniform magnetic field of same direction \& strength $-B\hat k$ as in region-I but it can be switched on \& off, after electrons have entered by some distance in region-II. [Charge of e$^-$ is $e$ (magnitude) \& mass $=m$]
\qfig{p1_phy_s2_para9_10_1}
\qn{9} Choose correct statement regarding possible path of electrons in region-II, when magnetic field is periodically changed.
\begin{qoptions}
\opt{A}{If magnetic field is switched on \& off in $\dfrac{\pi m}{eB}$ time interval then path of electron will be :\par\qfigi{p1_phy_s2_q09_1}}
\opt{B}{If magnetic field is switched on \& off in $\dfrac{\pi m}{2eB}$ time interval then path of electron will be:\par\qfigi{p1_phy_s2_q09_2}}
\opt{C}{If magnetic field is switched on and off in time interval of $\dfrac{2\pi m}{eB}$ time interval path of electron will be :\par\qfigi{p1_phy_s2_q09_3}}
\opt{D}{If magnetic field is switched on and off in time interval of $\dfrac{2\pi m}{3eB}$ then path of electron will be:\par\qfigi{p1_phy_s2_q09_4}}
\end{qoptions}
\qn{10} If magnetic field is switched on \& off in $\dfrac{2\pi m}{eB}$ time interval then average speed with which electron will appear to move along $x$-axis, during long interval of time is :
\begin{qoptions}
\opt{A}{$\dfrac{E}{2B}$}
\opt{B}{$\dfrac EB\left(\dfrac{\pi-2}{\pi}\right)$}
\opt{C}{$\dfrac EB\left(\dfrac{\pi-2}{2}\right)$}
\opt{D}{electron will not move forward as it is stuck in a loop.}
\end{qoptions}
\sect{Section-II}
\qn{1} If peak voltage of source is 100V, then peak value of I is $\dfrac{\sqrt{100x}}{5}$, then value of $x$ is .......
\qfig{p1_phy_s3_q01_1}
\qn{2} If magnetic energy of indutor becomes 40\% of initial value in 4 seconds then at what time (in sec) this energy becomes 16\% of initial value, if switch 's' is closed at $t=4$ sec.
\qfig{p1_phy_s3_q02_1}
\qn{3} A particle having mass $m$ and charge $+q$ is released from rest in uniform electric and magnetic field perpendicular to each other then maximum speed of particle (m/s) during motion will be : ($E=232$ N/c, $B=5$T, $q=2$mc, $m=1$gm)
\qn{4} A small current carrying loop is slowly turned in uniform magnetic field $\vec B=\left(1.2\,\hat i+2.5\hat j+3\hat k\right)$ T such that initial and final magnetic moment of loop is $\left(4\hat i+3\hat j-12\hat k\right)$ J/T and $\left(3\hat i+12\hat j-8\hat k\right)$ J/T respectively work done in turning will be (in J) :
\qn{5} Initially key $k_1$ is closed and $k_2$ is open for long time, then suddenly $k_1$ is opened and $k_2$ is closed, simultaneously. Rate of change (in A/s) in current in inductor just after this change is....
\qfig{p1_phy_s3_q05_1}
\qn{6} Magnetic field $\vec B$ is present in region between circles having radii R and 2R. $B=(7t+4)$T, $t$ in sec. How much work (in J) has to be done to displace a particle having charge $q=200$ mC from 1 to 2 as shown in figure.
\qfig{p1_phy_s3_q06_1}
\qn{7} Two conducting rods AB and CD have resistance $3\Omega$ and $6\Omega$ respectively. Find out current (in A) in resistor R. [R $=2\Omega$, $B=2.4$T, $v=5$ m/s, $\ell=2$m]
\qfig{p1_phy_s3_q07_1}
\qn{8} Quality factor of a series LCR circuit is 0.6. If definition of quality factor is changed such that now limiting frequency of bandwidth is taken upto $\dfrac13$rd of maximum power rather than half of maximum power (standard definition), then what will be quality factor of given circuit ?
\partsec{2}{CHEMISTRY}
\sect{Section-I(i)}
\qn{1} Which of the following reaction(s) give a product which on reaction with solution of AgNO$_3$ / NH$_4$OH gives white ppt.
\begin{qoptions}
\opt{A}{$\mathrm{CaC_2}\xrightarrow{\;\mathrm{H_2O}\;}$}
\opt{B}{$\mathrm{CHCl_3}\xrightarrow[\Delta]{\mathrm{Ag}}$}
\opt{C}{$\mathrm{Mg_2C_3}\xrightarrow{\;\mathrm{H_2O}\;}$}
\opt{D}{$\mathrm{Al_4C_3}\xrightarrow{\;\mathrm{H_2O}\;}$}
\end{qoptions}
\qn{2} Which among the following reaction(s) would result formation of \qfigi{p1_chem_s1_q02_1} as a major product ?
\begin{qoptions}
\opt{A}{\qfigi{p1_chem_s1_q02_A}$\xrightarrow{\text{\footnotesize P.C.C. (Pyridinium chloro chromate)}}$}
\opt{B}{\qfigi{p1_chem_s1_q02_B}$\xrightarrow[1:1]{\text{Et--MgCl}}$}
\opt{C}{\qfigi{p1_chem_s1_q02_C}$\xrightarrow[\text{(2) H}_2\text{O}]{\text{(1) }\qfigi{p1_chem_s1_q02_E}}$}
\opt{D}{\qfigi{p1_chem_s1_q02_D}$\xrightarrow[\text{\footnotesize (2) P.D.C. (Pyridinium dichromate)}]{\text{\footnotesize (1) dil. H}_2\text{SO}_4}$}
\end{qoptions}
\qn{3} Which of the following abbreviation is/are represent bidentate ligands (maximum denticity) ?
\begin{qoptions}
\opt{A}{dien}
\opt{B}{dipy}
\opt{C}{pn}
\opt{D}{gly$^{-1}$}
\end{qoptions}
\qn{4} Which complex compound/s do/ does not follow the effective Atomic Number Rule ?
\begin{qoptions}
\opt{A}{$[\mathrm{Fe}(\pi\text{-}\mathrm{C_5H_5})_2]$}
\opt{B}{$[\mathrm{Fe(H_2O)_5NO}]\,\mathrm{SO_4}$}
\opt{C}{$\mathrm{Na_2}\,[\mathrm{Fe(CN)_5\,NO}]$}
\opt{D}{$\mathrm{K_3}\,[\mathrm{Fe\,(CN)_6}]$}
\end{qoptions}
\qn{5} At 25$^\circ$C, conductivity of aq. solution of 0.01 M Na$_2$SO$_4$ is found to be 0.18 S/m. Given that limiting molar conductivity of $\mathrm{SO_4^{2-}}$ ion is 80 S cm$^2$ mol$^{-1}$, then select the correct options.
\begin{qoptions}
\opt{A}{Molar conductance of 0.01 M Na$_2$SO$_4$ (aq.) solution is $1.8\times10^4$ S cm$^2$ mol$^{-1}$}
\opt{B}{Equivalent conductance of 0.01 M Na$_2$SO$_4$ solution is approximately 90 S cm$^2$ eq$^{-1}$.}
\opt{C}{Limiting molar conductance of Na$^+$ is approximately 100 S cm$^2$ mol$^{-1}$}
\opt{D}{Limiting equivalent conductance of Na$^+$ is approximately 50 S cm$^2$ eq$^{-1}$}
\end{qoptions}
\qn{6} 70 mL of 0.1 M HCl (aq.) is titrated against 0.2 M K$_2$Cr$_2$O$_7$ (aq.) solution to form CrCl$_3$, Cl$_2$ and KCl. Select the correct options :
\begin{qoptions}
\opt{A}{Required volume of K$_2$Cr$_2$O$_7$ to titrate HCl completely will be $\dfrac{35}{6}$mL}
\opt{B}{Moles of Cl$_2$ formed during reaction is $1.5\times10^{-3}$}
\opt{C}{Equivalent mass of CrCl$_3$ is $\dfrac M6$ (M = formula mass of CrCl$_3$)}
\opt{D}{Equivalent mass of Cl$_2$ is $\dfrac M2$ (M = Molar mass of Cl$_2$)}
\end{qoptions}
\sect{Section-I(ii)}
\paraq{7}{8}
Calculation of percentage of carbon and hydrogen it is determined by Liebig's method. It is based on the fact that an organic compound on heating with dry copper oxide gives CO$_2$ and water. Water is absorbed in conc. H$_2$SO$_4$ while CO$_2$ is absorbed by KOH solution.

Organic compound $+\ \mathrm{O_2}\xrightarrow{\;\mathrm{CuO}\;}\mathrm{CO_2}+\mathrm{H_2O}$

Thus by knowing the amount of CO$_2$ and H$_2$O produced from a known weight of organic compound the \% of carbon and hydrogen can be calculated.

E.M.F. of a Galvanic all cell given by a thermodynamic equation known as Nernst equation.
\[E=E^\circ-\dfrac{RT}{nF}\,\ell n\,Q^{\circ}_{pc}\]
\qn{7} 6 gm of a weak organic acid (HA) on heating with dry copper oxide give 8.8 gm CO$_2$ and 3.6 gm water. Molecular formula of the acid can be
\begin{qoptions}
\opt{A}{HCOOH}
\opt{B}{H$_2$C$_2$O$_4$}
\opt{C}{CH$_3$COOH}
\opt{D}{C$_2$H$_5$COOH}
\end{qoptions}
\qn{8} 6 gm of above acid dissolved in 1 litre of water and this solution is used to make a hydrogen electrode that is clubed with standard calomel electrode at T K

$\mathrm{Pt|H_2(1bar)|HA||}$ Standard calomel electrode

e.m.f. of this cell is found to be 0.46 V value of K$_a$ for weak acid HA will be

(Given that $\dfrac{2.303RT}{F}=0.06$, $E^\circ_{\mathrm{Cl^-|Hg_2Cl_2|Hg|Pt}}=0.28$ V)
\begin{qoptions}
\opt{A}{$2\times10^{-5}$}
\opt{B}{$10^{-6}$}
\opt{C}{$10^{-4}$}
\opt{D}{$10^{-5}$}
\end{qoptions}
\paraq{9}{10}
For given reaction sequence molecular formula for compound 'U' is C$_7$H$_6$O$_2$ \& "P" gives negative Fehling test.
\qfig{p1_chem_s2_para9_10_1}
\qn{9} Compound which is not a hydrocarbon
\begin{qoptions}
\opt{A}{Compound (W)}
\opt{B}{Compound (R)}
\opt{C}{Compound (T)}
\opt{D}{Compound (V)}
\end{qoptions}
\qn{10} $\mathrm{Q}\xrightarrow{\;\mathrm{NaOH}+\mathrm{I_2}\;}\mathrm{X}+\mathrm{Y}$ (yellow ppt.) :

$\mathrm{P}\xrightarrow[\mathrm{CH_3COO^{\ominus},\,\Delta}]{(\mathrm{CH_3CO})_2\mathrm{O}}\mathrm{Z}$

identify the correct statement
\begin{qoptions}
\opt{A}{Compound Z is cinnamaldehyde}
\opt{B}{$\mathrm{Y}\xrightarrow[\text{Powder, }\Delta]{\mathrm{Ag}}\mathrm{W}$}
\opt{C}{$\mathrm{X}\xrightarrow[\mathrm{CaO,\,\Delta}]{\mathrm{NaOH}}\mathrm{CH_3-CH_3}$}
\opt{D}{DOU of compound Z is 5.}
\end{qoptions}
\sect{Section-II}
\qn{1} What is the sum of point numbers that are applicable to the process of calcination of an ore in the concentration step ? (For example 1, 2 and 4$^{\text{th}}$ statement are applicable then answer of question is $1+2+4=7$)

(1) Strong heating

(2) Presence of oxygen

(3) Expulsion of volatile organic impurities

(4) Formation of volatile oxides of nonmetallic impurities.

(5) Removal of moisture

(6) Decomposition of carbonates

(7) Fusion of ore to form molten mass

(8) Resulting substance becomes porous.
\qn{2} For the complex $[\mathrm{NiF_6}]^{-2}$ ion the magnitude of Crystal Field Stablization Energy in term of $\Delta_0$ is .............. (neglecting pairing energy)
\qn{3} Excess of cold very dilute solution of nitric acid reacts with 26.152 g of pure zinc to form zinc nitrate and x grams of a nitrogen containing compound. Find the value x.

(Atomic mass of H $=1$, N $=14$, O $=16$, Zn $=65.38$)
\qn{4} For the complex ion $[\mathrm{Co\,(H_2O)_2\,(OH)_2\,(C_2O_4)}]^{-}$ if

$x$ = Number of unpaired electrons

$y$ = Number of geometrical isomers

$z$ = Number of stereo isomers

Calculate $x+y+z$
\qn{5} How many of the given compounds shows +ve tollen's test :
\qfig{p1_chem_s3_q05_1}
(V) CH$_3$ --- CHO \hfill (VI) HCOOH\par
(VII) CH$_3$COOH \hfill (VIII) Benzophenone\par
(IX) Acetophenone \hfill (X) Pyruvaldehyde
\qn{6} $\mathrm{C_{14}H_{20}O}\ \mathrm{(X)}\xrightarrow[\mathrm{HCl}]{\mathrm{Zn/Hg}}$ \qfigi{p1_chem_s3_q06_1}

Find all possible X ?
\qn{7} During electrolysis of aqueous conctrated NaCl using Pt-electrodes, if 3 moles of water is electrolyted then how much charge (in Faradays) is required if current effeciency is 75\%
\qn{8} 6.2 g of a sample containing Na$_2$CO$_3$, NaHCO$_3$ \& non-volatile inert impurity on gentle heating loses 5\% of its weight due to formation of H$_2$O and CO$_2$ as in reaction $2\mathrm{NaHCO_3}\rightarrow\mathrm{Na_2CO_3}+\mathrm{H_2O}+\mathrm{CO_2}$.

Residue is dissolved in water and formed 100 ml solution and its 10 ml portion requires 7.5 ml of 0.2 M aqueous solution of BaCl$_2$ for complete precipitation of carbonates.

Determine weight (in gram) of Na$_2$CO$_3$ in the original sample
\partsec{3}{MATHEMATICS}
\sect{Section-I(i)}
\qn{1} If P is any point on the ellipse $\dfrac{x^2}{a^2}+\dfrac{y^2}{b^2}=1$ and S is the focus of the ellipse and $\dfrac xy$ (where $x$ \& $y$ are co-prime natural numbers) is ratio of the area of the ellipse and the area enclosed by the locus of mid-point of PS then which of the following is/are true
\begin{qoptions}
\opt{A}{$x+y=5$}
\opt{B}{$|y-x|=3$}
\opt{C}{$x+y=10$}
\opt{D}{$x^2+y^2=17$}
\end{qoptions}
\qn{2} If the value of $\dfrac{\log_{81}729+\log_{729}81}{\log_{\sqrt{\frac94}}\sqrt{\frac{81}{16}}+\log_{\sqrt{\frac{81}{16}}}\sqrt{\frac94}}=\dfrac ab$, then the value of $(b-a)$ is equal to, (where $a$ \& $b$ are coprime numbers) -
\begin{qoptions}
\opt{A}{$b-a=2$}
\opt{B}{$b+a=28$}
\opt{C}{$2b-a=17$}
\opt{D}{$2a-b=11$}
\end{qoptions}
\qn{3} A variable line intersects co-ordinate axes at two points A and B. If the area of traingle OAB (where 'O' is origin) is 4 square units, then the locus of middle point of AB is -
\begin{qoptions}
\opt{A}{$xy=2$}
\opt{B}{$xy=-2$}
\opt{C}{$x^2+y^2=4$}
\opt{D}{$x^2-y^2=4$}
\end{qoptions}
\qn{4} A line $kx-y=0$ intersects the curve $y=1-\sqrt{x-1}$ at -
\begin{qoptions}
\opt{A}{no point for $k\in\left(\dfrac{1-\sqrt2}{2},1\right)$}
\opt{B}{no point for $k\in\left(-\infty,\dfrac{1-\sqrt2}{2}\right)\cup(1,\infty)$}
\opt{C}{two points for $k\in\left(\dfrac{1-\sqrt2}{2},0\right)$}
\opt{D}{one point for $k\in[0,\infty)$}
\end{qoptions}
\qn{5} If the maximum and minimum value of the expression $y=\dfrac{\sqrt2\sin x}{3-\cos x}$ are $m$ and $n$ respectively, then :
\begin{qoptions}
\opt{A}{$\dfrac mn=-1$}
\opt{B}{$3m+n=1$}
\opt{C}{$m^2+n^2=\dfrac14$}
\opt{D}{$m-n=1$}
\end{qoptions}
\qn{6} On an ellipse $\dfrac{x^2}{64}+\dfrac{y^2}{9}=1$, tangents drawn at $P_1,P_2,\dots P_8$ intersects the major axis at $T_1,T_2,\dots T_8$ respectively. If the value of $\displaystyle\sum_{i=1}^{8}\dfrac{(\text{area of }\Delta P_iT_iF_1)(\text{area of }\Delta P_iT_iF_2)}{(P_iT_i)^2}=N$, where $F_1$ and $F_2$ are the foci of ellipse, then
\begin{qoptions}
\opt{A}{$N$ is perfect square}
\opt{B}{$N$ is odd number}
\opt{C}{$N$ is prime number}
\opt{D}{$N\le25$}
\end{qoptions}
\sect{Section-I(ii)}
\paraq{7}{8}
The tangents from the point P(2, 4) to a parabola touching it at A(4, 6) \& B(4, 2). The normals to parabola at A \& B intersects at point D. Point P, mid-point of AB \& point D are collinear.

Let C $\equiv x^2+y^2=r^2$.

On the basis of above information answer the following :
\qn{7} If there exists exactly one point in $x$-$y$ plane from where it is possible to draw two perpendicular tangents to parabola as well as circle 'C', then radius of circle 'C' is :-
\begin{qoptions}
\opt{A}{4}
\opt{B}{2}
\opt{C}{$\sqrt2$}
\opt{D}{none of these}
\end{qoptions}
\qn{8} Let $r=1$ for C, and $(x_1,y_1)$ be a point on the parabola which is at a minimum distance from C, then the value of $(x_1+2)(y_1-4)$ is :-
\begin{qoptions}
\opt{A}{8}
\opt{B}{$-8$}
\opt{C}{4}
\opt{D}{$-4$}
\end{qoptions}
\paraq{9}{10}
Let $\sin\theta-\sin\left(\dfrac\theta2\right)=a$ \& $\cos\theta+\cos\dfrac\theta2=b$.

On the basis of above information, answer the following questions :
\qn{9} The value of $a^2+b^2$ when $\theta=80^\circ$ is -
\begin{qoptions}
\opt{A}{1}
\opt{B}{0}
\opt{C}{$\dfrac12$}
\opt{D}{$-1$}
\end{qoptions}
\qn{10} The value of $b^2-a^2$ is -
\begin{qoptions}
\opt{A}{$4\cos^2\dfrac{3\theta}{4}\cos\dfrac\theta2$}
\opt{B}{$4\cos^2\dfrac\theta2\cos\dfrac{3\theta}{4}$}
\opt{C}{$4\cos\dfrac\theta2\cos\dfrac{3\theta}{4}$}
\opt{D}{$4\cos^2\dfrac\theta2\cos^2\dfrac{3\theta}{4}$}
\end{qoptions}
\sect{Section-II}
\qn{1} If $x=\log_3\left[\dfrac{\sqrt{30+\sqrt{30+\sqrt{30+\dots\infty}}}}{\sqrt{2+\sqrt{2+\sqrt{2+\dots\infty}}}}\right]$, find the value of $x$ ?
\qn{2} Let $C_1$ be the graph of $xy=1$ and the reflection of $C_1$ in the line $y=2x$ is $C_2$. If the equation of $C_2$ is expressed as $12x^2+bxy+cy^2+d=0$ then (bc) is equal to
\qn{3} Consider an ellipse with length of major axis as 12 and length of minor axis as 6. Let A, B, C, D are points on ellipse such that tangents to the ellipse at these points form a square, then find area of square -
\qn{4} In a quadrilateral ABCD if $\sin\dfrac{A+B}{2}\cos\dfrac{A-B}{2}+\sin\dfrac{C+D}{2}\cdot\cos\dfrac{C-D}{2}=2$, then $\sum\cos\dfrac A2\cdot\cos\dfrac B2$ is equal to
\qn{5} If $x<0$, \& $x^2-x-1=0$, then the value of $\left(\log_{10}(x^4+1)(x^8+1)(x^{16}+1)-14\log_{10}(-x)\right)$ is $\log_{10}k$. Find k
\qn{6} Let P and Q be the end points of the latus rectum and R$(-5\cos45^\circ,-5\sin45^\circ)$ be the vertex of the parabola. Circle $x^2+y^2-25=0$ is circumcircle of $\Delta$PQR. If point $(a,b)$ is intersection point of tangents at the points P and Q to the parabola, then the value of $(a^2+b^2)$ is
\qn{7} If variable chord of hyperbola $x^2-y^2=9$ touches the parabola $y^2=12x$ and locus of middle point of these chords is $x^3+\lambda_1xy^2+\lambda_2y^2=0$, then $(\lambda_1+\lambda_2)$ equals
\qn{8} If P$(\alpha,\beta)$ is mid point of the chord intercepted on the line $9x-4y=13$ by the ellipse $\dfrac{x^2}{4}+\dfrac{y^2}{9}=1$ then find CP$^2$, where C is center of ellipse
\papersec{2}
\partsec{1}{PHYSICS}
\sect{Section-I}
\qn{1} Net magnetic field at the centre of the circle O due to current carrying loop shown by solid line is
\qfig{p2_phy_s1_q01_1}
\begin{qoptions}
\opt{A}{Zero}
\opt{B}{Perpendicular to sheet inwards}
\opt{C}{Perpendicular to sheet outwards}
\opt{D}{Cannot be determined with given information}
\end{qoptions}
\qn{2} Consider an amperian loop PQRST (QS is diameter of semicircle QRS) around an infinitely long current carrying wire at O(current comes out of plane). The value of line integral $\displaystyle\oint_{\text{QPTS}}\vec B\bullet d\vec\ell$ (traversed anti clockwise) is:
\qfig{p2_phy_s1_q02_1}
\begin{qoptions}
\opt{A}{$\dfrac{\mu_0 i}{2}$}
\opt{B}{$>\dfrac{\mu_0 i}{2}$}
\opt{C}{$<\dfrac{\mu_0 i}{2}$}
\opt{D}{$\dfrac{4\mu_0 i}{\pi}$}
\end{qoptions}
\qn{3} For the given circuit, switch is closed at $t=0$. Choose the incorrect statement.
\qfig{p2_phy_s1_q03_1}
\begin{qoptions}
\opt{A}{For $t>0$, current through $10\ \Omega$ resistor decreases with time and finally becomes zero.}
\opt{B}{Current through inductor at steady state is 2.4 A.}
\opt{C}{Current through battery just after closing the switch is 0.8 A.}
\opt{D}{For $t>0$, current through $5\ \Omega$ resistor decreases with time and finally reaches a constant value.}
\end{qoptions}
\qn{4} The given structure LMNO is a conducting semicircular ring MNO (radius $=6$cm) joined to a conducting straight wire LM (5cm). It is hinged at `O' and rotated about it with constant angular speed $\omega$ in a region with uniform magnetic field perpendicular to plane of structure. The ratio of induced EMF in the curved quarter segment MN to that induced in straight segment LM is \underline{\hspace{2.5em}}.
\qfig{p2_phy_s1_q04_1}
\begin{qoptions}
\opt{A}{$\dfrac{144}{25}$}
\opt{B}{$\dfrac{72}{25}$}
\opt{C}{$\dfrac{36}{25}$}
\opt{D}{$\dfrac{25}{36}$}
\end{qoptions}
\qn{5} In the series RLC circuit, the voltage output of AC source is given by $V_s=30\cos(50\pi t)$ Volts. The value of instantaneous potential difference between points A and B, when the applied voltage has instantaneous value 18V is \underline{\hspace{2.5em}}.
\qfig{p2_phy_s1_q05_1}
\begin{qoptions}
\opt{A}{18V}
\opt{B}{0V}
\opt{C}{12V}
\opt{D}{24V}
\end{qoptions}
\qn{6} A negatively charged particle $(m,-q,v)$ enters a region of uniform magnetic field (B) at an angle of $30^\circ$ with the boundary as shown in figure. The magnitude of average magnetic force on charged particle averaged during the time interval it remains in magnetic field is:
\qfig{p2_phy_s1_q06_1}
\begin{qoptions}
\opt{A}{$\dfrac{3qBV}{5\pi}$}
\opt{B}{$qBV$}
\opt{C}{$\dfrac{6qBV}{\pi}$}
\opt{D}{$\dfrac{qBV}{\pi}$}
\end{qoptions}
\qn{7} In the given LC circuit, charge on the right plate `q' of capacitor and current through inductor are mentioned for time $t=0$. If natural frequency of oscillation is $f=\dfrac{10^4}{8\pi}$Hz, write expression for $q(t)$ i.e. charge on right plate of capacitor as a function of time (for $t\ge0$).
\qfig{p2_phy_s1_q07_1}
\begin{qoptions}
\opt{A}{$q(t)=1\,\text{mC}\,\sin(2500\,t-37^\circ)$}
\opt{B}{$q(t)=1\,\text{mC}\,\sin(2500\,t+37^\circ)$}
\opt{C}{$q(t)=1\,\text{mC}\,\cos(2500\,t+53^\circ)$}
\opt{D}{$q(t)=1\,\text{mC}\,\cos(2500\,t-53^\circ)$}
\end{qoptions}
\qn{8} Choose the incorrect statement regarding transformers:
\begin{qoptions}
\opt{A}{In step-up transformers, the primary coil is made from thick wire so that it can sustain high current.}
\opt{B}{Transformers do not work on DC voltage}
\opt{C}{In transformers, EMF induced in secondary coil is $180^\circ$ out of phase with EMF applied across primary coil.}
\opt{D}{In transformers, steel is preferred over soft-iron as cores inside current carrying coils}
\end{qoptions}
\qn{9} A long solenoid is fabricated by closely winding an insulated wire of cross-sectional diameter 0.5 mm, over a cylindrical long rod of material of relative permeability $\mu_r=200$. Successive turns nearly touch each other. Find magnetic field at centre of solenoid if wire carries current of 2A.
\qfig{p2_phy_s1_q09_1}
\begin{qoptions}
\opt{A}{$\dfrac{8\pi}{25}$ Tesla}
\opt{B}{$\dfrac{4\pi}{25}$ Tesla}
\opt{C}{$\dfrac{8\pi}{25}\times10^{-3}$ Tesla}
\opt{D}{$\dfrac{4\pi}{25}\times10^{-3}$ Tesla}
\end{qoptions}
\qn{10} A long solenoid of radius R carries a time dependent current $i(t)=5t\,(3-2t)$ Ampere. A square shaped conducting loop of side length 4R is placed co-axially near centre of solenoid. For time interval $0\le t\le2$ sec, the magnetic flux $(\phi_s)$, induced current $(I_s)$ and induced EMF $(E_s)$ in the square loop varies as:
\qfig{p2_phy_s1_q10_1}
\begin{qoptions}
\opt{A}{$E_s$ is zero at $t=0$ and $t=1.5$ sec. $I_s$ is maximum at $t=0.75$ sec.}
\opt{B}{$\phi_s$ is maximum at $t=0.75$ sec. $I_s$ reverses its direction at $t=0.75$ sec.}
\opt{C}{$\phi_s$ is zero at $t=1.5$ sec. $I_s$ reverses its direction at $t=0.75$ sec.}
\opt{D}{For the interval $0.75\le t\le2$ sec, $I_s$ first decreases and then increases.}
\end{qoptions}
\qn{11} A voltage source $V=200\sqrt2\sin(\omega t+\theta)$ connected in series LCR circuit having resistance, inductance and capacitance as $100\ \Omega$, 1H and $9\ \mu$f respectively. Then max power consumed will be
\begin{qoptions}
\opt{A}{400 w}
\opt{B}{800 w}
\opt{C}{200 w}
\opt{D}{100 w}
\end{qoptions}
\qn{12} A barmagnet with pole strength `m' and magnetic moment `M' is divided in two equal halves by two method, as shown below in given figure.
\qfig{p2_phy_s1_q12_1}
Then new pole strength in each half in a and b respectively is:
\begin{qoptions}
\opt{A}{$\dfrac m2,\dfrac m2$}
\opt{B}{$m,m$}
\opt{C}{$\dfrac m2,m$}
\opt{D}{$m,\dfrac m2$}
\end{qoptions}
\qn{13} Choose the correct statement:
\begin{qoptions}
\opt{A}{Earth magnetic field is due to large magnet placed inside it.}
\opt{B}{Angle between geographic meridian and magnetic meridian is defined as angle of dip.}
\opt{C}{Angle made by net magnetic field with horizontal at a point is defined as angle of dip.}
\opt{D}{The angle made by magnetic needle with horizontal in dip circle in a plane perpendicular to plane where needle is horizontal is there angle of dip.}
\end{qoptions}
\qn{14} A damped LC oscillation system has to be designed for critical damping. If value of inductance is 4H and capacitance is $1\mu$F, then resistance required should be
\begin{qoptions}
\opt{A}{$2000\ \Omega$}
\opt{B}{$2\ \Omega$}
\opt{C}{$1000\ \Omega$}
\opt{D}{$100\ \Omega$}
\end{qoptions}
\qn{15} A solenoid has an inductance of 10 Henry and a capacitor is connected to form a LC oscillator. The charge on capacitor varies as $Q=Q_0\sin\left(\omega t+\dfrac\pi3\right)$. If the time period of one oscillation is 24 sec, then minimum time after which energy stored in inductor is thrice of energy stored in capacitor is:
\begin{qoptions}
\opt{A}{12 sec}
\opt{B}{8 sec}
\opt{C}{4 sec}
\opt{D}{6 sec}
\end{qoptions}
\qn{16} A pure inductor is connected across AC source. Then which graph shows correct variation of peak current $i_0$ Vs frequency variation ?
\begin{qoptions}
\opt{A}{\qfigi{p2_phy_s1_q16_A}}
\opt{B}{\qfigi{p2_phy_s1_q16_B}}
\opt{C}{\qfigi{p2_phy_s1_q16_C}}
\opt{D}{\qfigi{p2_phy_s1_q16_D}}
\end{qoptions}
\qn{17} An alternating source of 200 V (RMS) is applied in series LCR circuit. At resonance potential difference across inductor and capacitor is 300 volts. Then potential difference across (RMS value) resistance is:
\begin{qoptions}
\opt{A}{300}
\opt{B}{$200\sqrt2$}
\opt{C}{200}
\opt{D}{Zero}
\end{qoptions}
\qn{18} Given $2X_L=X_C=R$. Maximum current through source is :
\qfig{p2_phy_s1_q18_1}
\begin{qoptions}
\opt{A}{$\dfrac{V_o}{R}$}
\opt{B}{$\dfrac{3V_o}{R}$}
\opt{C}{$\dfrac{2V_o}{R}$}
\opt{D}{0}
\end{qoptions}
\qn{19} A solid hemisphere having charge `Q' uniformly distributed on it. It is rotating with angular velocity `$\omega$' as shown, then magnetic moment is given by expression:
\qfig{p2_phy_s1_q19_1}
\begin{qoptions}
\opt{A}{$Q\omega R^2$}
\opt{B}{$\dfrac{2Q\omega R^2}{5}$}
\opt{C}{$\dfrac{Q\omega R^2}{2}$}
\opt{D}{$\dfrac{Q\omega R^2}{5}$}
\end{qoptions}
\qn{20} Small loop of side length `b' is placed at the centre of a big square loop, $(a\gg b)$ then mutual inductance of system is:
\qfig{p2_phy_s1_q20_A}
\begin{qoptions}
\opt{A}{$\dfrac{\mu_0\,2\sqrt2\,b^2}{\pi a}$}
\opt{B}{$\dfrac{\mu_0\,4b^2}{\pi a}$}
\opt{C}{$\dfrac{\mu_0\,2b^2}{2\pi a}$}
\opt{D}{$\dfrac{\mu_0\,b^2}{\sqrt2\pi a}$}
\end{qoptions}
\sect{Section-II}
\qn{1} In the given AC circuit, the average power dissipated in Watts is \underline{\hspace{2.5em}}.
\qfig{p2_phy_s2_q01_1}
\qn{2} A non-uniform magnetic field exists in a cylindrical region of radius (R) with centre at `O'. Magnetic field increases at the rate $\dfrac{dB}{dt}$. A proton and a deuterium ion are simultaneously released from points `P' and `S' respectively. If $\vec a_P$ and $\vec a_D$ are accelerations of proton and deuterium ion just after release then $\dfrac{\left|\vec a_P+\vec a_D\right|}{\left|\vec a_P\right|}=$? (Given $r_1=\dfrac R4$, $r_2=2R$ and OP is perpendicular to OS)
\qfig{p2_phy_s2_q02_1}
\qn{3} A non-uniform magnetic field B exists in a cylindrical region of radius 0.5m as shown. A uniform conducting wire of length 8m and total resistance $4\Omega$ is bent into a square frame and is placed with one side along the diameter of cylinder as shown. If $\dfrac{dB}{dt}=1$ Tesla/sec, find current induced in frame (in mA) (Rounded off to 2$^{\text{nd}}$ decimal place)\underline{\hspace{2.5em}}. (Take : $\pi=3.14$)
\qfig{p2_phy_s2_q03_1}
\qn{4} A 0.5 m long coil of wire having 20 turns carries a current of 1 A. If an alloy bar whose B -- H relation varies as shown in graph is inserted inside the coil, find the net magnetic flux density (in Tesla) inside the coil.
\qfig{p2_phy_s2_q04_1}
\qn{5} A thin long solenoid carries a current of 8A and has 400 turns per meter length of solenoid. For the equivalent bar magnet of same size and magnetization vector $\vec M$ (magnetic moment/volume), value of $\dfrac{\left|\vec M\right|}{100}$ in $(\text{Am}^{-1})$ is \underline{\hspace{2.5em}}
\qn{6} In a step down transformer input voltage is 20000 volts and number of turns in primary are 4000. The supply voltage at which electricity is being distributed is 200 volts. Then number of turns in secondary required considering only 80\% of flux linkage are $8x$. Find $x$.
\qn{7} In a vertical plane P$_1$ making angle $30^\circ$ with magnetic meridian apparent angle of dip is $60^\circ$. The apparent angle of dip in vertical plane P$_2$ which is perpendicular to P$_1$ is $\tan^{-1}(n)$. Find $n$.
\qn{8} If current amplitude is 5A and angular frequency of source is 100 rad/sec and resistance is $100\ \Omega$, inductance is 4 H and capacitance is $20\ \mu$F. The peak value of source voltage is $250\sqrt n$. Find $n$.
\qn{9} A rod of length 3m is hinged at a point 1m from one end is being rotated with angular velocity $\omega=2$ rad/sec as shown in uniform magnetic field of 2 Tesla pointing inward and perpendicular to plane of rotation of rod. Then $(V_P-V_Q)$ in volts is:
\qfig{p2_phy_s2_q09_1}
\qn{10} For given circuit frequency of oscillation is $F=\dfrac{1}{2\pi\sqrt{2x}}$ kHz. Find $n$.
\qfig{p2_phy_s2_q10_1}
\partsec{2}{CHEMISTRY}
\sect{Section-I}
\qn{1} Incorrect order of CO bond length :
\begin{qoptions}
\opt{A}{$[\mathrm{M(CO)_3PF_3}]<[\mathrm{M(CO)_3PCl_3}]$}
\opt{B}{$[\mathrm{M(CO)_3PF_3}]>[\mathrm{M(CO)_3PCl_3}]$}
\opt{C}{$[\mathrm{M(CO)_3PF_3}]<[\mathrm{M(CO)_3PMe_3}]$}
\opt{D}{$[\mathrm{M(CO)_3PCl_3}]<[\mathrm{M(CO)_3PMe_3}]$}
\end{qoptions}
\qn{2} Which of the following is not optically active ?
\begin{qoptions}
\opt{A}{$[\mathrm{Co(en)_3}]^{3+}$}
\opt{B}{$[\mathrm{Cr(ox)_3}]^{3-}$}
\opt{C}{cis-$[\mathrm{CoCl_2(en)_2}]^+$}
\opt{D}{trans-$[\mathrm{CoCl_2(en)_2}]^+$}
\end{qoptions}
\qn{3} Which of the following are incorrectly matched ?
\begin{qoptions}
\opt{A}{Casseterite -- SnO$_2$}
\opt{B}{Malachite -- CuCO$_3$.Cu(OH)$_2$}
\opt{C}{Azurite -- Cu$_2$O}
\opt{D}{Copper glance -- Cu$_2$S}
\end{qoptions}
\qn{4} In the extraction of nickel by Mond process, the metal is obtained by :
\begin{qoptions}
\opt{A}{Electrochemical reduction}
\opt{B}{Thermal decomposition}
\opt{C}{Chemical reduction by aluminium}
\opt{D}{Reduction by carbon}
\end{qoptions}
\qn{5} The no. of $\sigma$ and $\pi$ - bonds in Fe$_2$(CO)$_9$ respectively are :
\begin{qoptions}
\opt{A}{$22\sigma$ and $15\pi$}
\opt{B}{$23\sigma$ and $15\pi$}
\opt{C}{$22\sigma$ and $16\pi$}
\opt{D}{$15\sigma$ and $8\pi$}
\end{qoptions}
\qn{6} Observe the following partial roasting process, $\mathrm{CuFeS_2}+\mathrm{O_2}\rightarrow\mathrm{A}+\mathrm{FeO}+\mathrm{SO_2}$ ; A is :-
\begin{qoptions}
\opt{A}{Cu$_2$O}
\opt{B}{CuO}
\opt{C}{Cu$_2$S}
\opt{D}{CuS}
\end{qoptions}
\qn{7} Which of the following statement is correct ?
\begin{qoptions}
\opt{A}{Magnesite is the ore of magnesium}
\opt{B}{A mineral containing ZnS and PbS together is separated in froth floatation process by adding NaCl as depressant}
\opt{C}{White bauxite is purified by Hall's process}
\opt{D}{Extraction of zinc from zinc blende is possible by self reduction}
\end{qoptions}
\qn{8} The reaction of permagnanate ions with hydrogen peroxide in an acidic solution gives Mn(II) salt and at the same time oxygen is released

(P) $2\mathrm{MnO_4^-}+\mathrm{H_2O_2}+6\mathrm{H^+}\rightarrow2\mathrm{Mn^{+2}}+3\mathrm{O_2}+4\mathrm{H_2O}$

(Q) $2\mathrm{MnO_4^-}+3\mathrm{H_2O_2}+6\mathrm{H^+}\rightarrow2\mathrm{Mn^{+2}}+4\mathrm{O_2}+6\mathrm{H_2O}$

(R) $2\mathrm{MnO_4^-}+5\mathrm{H_2O_2}+6\mathrm{H^+}\rightarrow2\mathrm{Mn^{+2}}+5\mathrm{O_2}+8\mathrm{H_2O}$

(S) $2\mathrm{MnO_4^-}+7\mathrm{H_2O_2}+6\mathrm{H^+}\rightarrow2\mathrm{Mn^{+2}}+3\mathrm{O_2}+10\mathrm{H_2O}$

Which of the above reaction is balanced correctly ?
\begin{qoptions}
\opt{A}{P}
\opt{B}{Q}
\opt{C}{R}
\opt{D}{S}
\end{qoptions}
\qn{9} For given half cell calculate half cell potential $\mathrm{Pt/H_2(1atm)\,|\,HA(0.1\ M)}$ given $K_a$ (HA) $=10^{-5}$.
\begin{qoptions}
\opt{A}{$-0.354$ V}
\opt{B}{$+0.177$ V}
\opt{C}{$-0.177$ V}
\opt{D}{None of these}
\end{qoptions}
\qn{10} In an organic compound of molar mass 108 gm mol$^{-1}$ C, H and N atoms are present in 9 : 1 : 3.5 by weight. Molecular formula can be :
\begin{qoptions}
\opt{A}{C$_6$H$_8$N$_2$}
\opt{B}{C$_7$H$_{10}$N}
\opt{C}{C$_5$H$_6$N$_3$}
\opt{D}{C$_4$H$_{18}$N$_3$}
\end{qoptions}
\qn{11} Calculate the temporary and permanent hardness of water sample having the following constituent per litre

Ca(HCO$_3$)$_2$ $=324$ mg \quad MgCl$_2$ $=190$ mg

NaCl $=1170$ mg \quad Mg(HCO$_3$)$_2$ $=146$ mg

CaSO$_4$ $=272$ mg
\begin{qoptions}
\opt{A}{400 ppm, 150 ppm}
\opt{B}{300 ppm, 400 ppm}
\opt{C}{200 ppm, 300 ppm}
\opt{D}{300 ppm, 300 ppm}
\end{qoptions}
\qn{12} Which of the following graph truly represents the titration of CH$_3$COOH solution against NaOH solution ? (Conductance V/s V$_{\mathrm{NaOH}}$)
\begin{qoptions}
\opt{A}{\qfigi{p2_chem_s1_q12_A}}
\opt{B}{\qfigi{p2_chem_s1_q12_B}}
\opt{C}{\qfigi{p2_chem_s1_q12_C}}
\opt{D}{\qfigi{p2_chem_s1_q12_D}}
\end{qoptions}
\qn{13} Molar conductance at infinite dilution of NH$_4$Cl, NaOH and NaCl are 129.8, 217.4 and 108.9 Ohm$^{-1}$ cm$^2$mol$^{-1}$, respectively. If the molar conductance of 0.01 N solution of NH$_4$OH is 28.596 Ohm$^{-1}$ cm$^2$mol$^{-1}$, then the degree of dissociation of NH$_4$OH at this temperature is:
\begin{qoptions}
\opt{A}{6.3 \%}
\opt{B}{12 \%}
\opt{C}{44.7 \%}
\opt{D}{0.12 \%}
\end{qoptions}
\qn{14} In the mixture of NaHCO$_3$ and Na$_2$CO$_3$, the volume of a given HCl required is $x$ ml with phenolphthalein indicator and further $y$ ml is required with methyl orange indicator in same titration. Hence, the volume of HCl for complete reaction of NaHCO$_3$ present in the original mixture is
\begin{qoptions}
\opt{A}{$y$}
\opt{B}{$x/2$}
\opt{C}{$(y-x)$}
\opt{D}{$2x$}
\end{qoptions}
\qn{15} Which of the following is not correct major product in the given reactions
\begin{qoptions}
\opt{A}{$\mathrm{CH_3{-}C{\equiv}C{-}CH_3}\xrightarrow{\text{Na in liq.NH}_3}$ \qfigi{p2_chem_s1_q15_1}}
\opt{B}{$\mathrm{CH_3{-}C{\equiv}CH}\xrightarrow[\text{(ii) }\bar{\mathrm O}\mathrm{H,H_2O,H_2O_2}]{\text{(i) B}_2\mathrm{H_6},\text{ THF}}\ \mathrm{CH_3{-}CH_2{-}CH{=}O}$}
\opt{C}{$\mathrm{CH_3{-}C{\equiv}CH}\xrightarrow{\text{Na in liq.NH}_3}\mathrm{CH_3{-}CH{=}CH_2}$}
\opt{D}{$\mathrm{CH{\equiv}C{-}CH{=}CH_2}\xrightarrow{\text{HCl(leq.)}}\ \begin{array}[t]{@{}l@{}}\mathrm{CH_2{=}C{-}CH{=}CH_2}\\[-2pt]\ \ \ \ \,|\\[-2pt]\ \ \ \ \mathrm{Cl}\end{array}$}
\end{qoptions}
\qn{16} Which of the following statement is not correct regarding the major product of following reaction:

\qfigi{p2_chem_s1_q16_1}$\xrightarrow{\text{Alc.KOH}}$ Major product
\begin{qoptions}
\opt{A}{Major product give optically inactive product mixture on reaction with Cl$_2$/CHCl$_3$}
\opt{B}{Total 6 isomers are possible having the molecular formula of major product}
\opt{C}{Oxidation of major product with cold alkaline KMnO$_4$ gives a MESO compound}
\opt{D}{It has more melting point as compared to one of it's isomer.}
\end{qoptions}
\qn{17} $\mathrm{CH_3-CH_2-CH{=}O+HCHO}\xrightarrow[\text{(excess)}]{\bar{\mathrm O}\mathrm H}$

Which of the following product mixture will contain major part ?
\begin{qoptions}
\opt{A}{$\begin{array}[t]{@{}l@{}}\mathrm{CH_3-CH-CH{=}O}\\[-2pt]\ \ \ \ \ \ \ \ \ |\\[-2pt]\ \ \ \ \ \ \ \ \mathrm{CH_2-OH}\end{array}$}
\opt{B}{$\begin{array}[t]{@{}l@{}}\ \ \ \ \ \ \ \ \ \mathrm{CH_2-OH}\\[-2pt]\ \ \ \ \ \ \ \ \ \ \ |\\[-2pt]\mathrm{CH_3-C-CHO+HCOO^-}\\[-2pt]\ \ \ \ \ \ \ \ \ \ \ |\\[-2pt]\ \ \ \ \ \ \ \ \ \mathrm{CH_2-OH}\end{array}$}
\opt{C}{$\begin{array}[t]{@{}l@{}}\ \ \ \ \ \ \ \ \ \mathrm{CH_2-OH}\\[-2pt]\ \ \ \ \ \ \ \ \ \ \ |\\[-2pt]\mathrm{CH_3-C-CH_2-OH+HCOO^-}\\[-2pt]\ \ \ \ \ \ \ \ \ \ \ |\\[-2pt]\ \ \ \ \ \ \ \ \ \mathrm{CH_2-OH}\end{array}$}
\opt{D}{$\begin{array}[t]{@{}l@{}}\ \ \ \ \ \ \ \ \ \mathrm{CH_2-OH}\\[-2pt]\ \ \ \ \ \ \ \ \ \ \ |\\[-2pt]\mathrm{CH_3-C-COO^-+CH_3OH}\\[-2pt]\ \ \ \ \ \ \ \ \ \ \ |\\[-2pt]\ \ \ \ \ \ \ \ \ \mathrm{CH_2-OH}\end{array}$}
\end{qoptions}
\qn{18} Which of the following reaction do not resemble with the major product given
\begin{qoptions}
\opt{A}{$\mathrm{R-C{\equiv}N}\xrightarrow[\text{(ii) H}_3\mathrm{O}^+]{\text{(i) SnCl}_2\text{+HCl}}\mathrm{R-CH{=}O}$}
\opt{B}{$\begin{array}[t]{@{}l@{}}\mathrm{R-C-Cl}\\[-2pt]\ \ \,\|\\[-2pt]\ \ \,\mathrm O\end{array}\xrightarrow[\text{(ii) H}_3\mathrm{O}^+]{\text{(i) H}_2/\text{Pd}/\mathrm{BaSO_4}}\mathrm{RCHO}$}
\opt{C}{$\begin{array}[t]{@{}l@{}}\mathrm{R-C-NH_2}\\[-2pt]\ \ \,\|\\[-2pt]\ \ \,\mathrm O\end{array}\xrightarrow{\mathrm{NaBH_4}}\mathrm{R-CH_2-NH_2}$}
\opt{D}{$\mathrm{R-CH{=}NH}\xrightarrow{\mathrm{NaBH_4}}\mathrm{R-CH_2-NH_2}$}
\end{qoptions}
\qn{19} Which reaction will not give alkyne as major product ?
\begin{qoptions}
\opt{A}{$\begin{array}[t]{@{}l@{}}\ \ \ \ \ \ \ \mathrm{Br\ \ Br}\\[-2pt]\ \ \ \ \ \ \ \ |\ \ \ \ \ |\\[-2pt]\mathrm{CH_3-C-C-CH_3}\\[-2pt]\ \ \ \ \ \ \ \ |\ \ \ \ \ |\\[-2pt]\ \ \ \ \ \ \ \mathrm{Br\ \ Br}\end{array}$$\xrightarrow[\Delta]{\text{Zn dust}}$}
\opt{B}{$\mathrm{CH_3-CHBr-CHBr-CH_3}\xrightarrow[\text{(ii) NaNH}_2]{\text{(i) Alc. KOH}}$}
\opt{C}{$\mathrm{BrCH_2-CH_2Br}\xrightarrow{\text{aquous KOH}}$}
\opt{D}{$\mathrm{CH_3-CHBr-CH_2Br}\xrightarrow[\text{(ii) NaNH}_2]{\text{(i) Alc. KOH}}$}
\end{qoptions}
\qn{20} Which of the following reaction of Alkene match with the correct major product ?
\begin{qoptions}
\opt{A}{\qfigi{p2_chem_s1_q20_A1}$\xrightarrow[\text{(ii) NaBH}_4,\bar{\mathrm O}\mathrm H]{\text{(i) Hg(OAC)}_2,\text{THF,H}_2\mathrm O}$ \qfigi{p2_chem_s1_q20_A2}\ Optically pure product}
\opt{B}{\qfigi{p2_chem_s1_q20_B1}$\xrightarrow[\text{(ii) CH}_3\text{COOH,H}_2\mathrm O]{\text{(i) B}_2\mathrm H_6,\text{THF}}$ \qfigi{p2_chem_s1_q20_B2}}
\opt{C}{\qfigi{p2_chem_s1_q20_C1}$\xrightarrow[\text{(ii) H}_3\mathrm O^+]{\text{(i) MCPBA}}$ \qfigi{p2_chem_s1_q20_C2}}
\opt{D}{\qfigi{p2_chem_s1_q20_D1}$\xrightarrow[\text{(ii) H}_2\mathrm O_2]{\text{(i) O}_3}2\mathrm{CH_3CHO}$}
\end{qoptions}
\sect{Section-II}
\qn{1} How many of the following molecules have large hydration constant ?
\qfig{p2_chem_s2_q01_1}
\qn{2} Consider the following reactions :

(1) \qfigi{p2_chem_s2_q02_a}$\xrightarrow{\mathrm{H^+}/\Delta}$ P (Major product)

(2) \qfigi{p2_chem_s2_q02_b}$\xrightarrow{\mathrm{H^+}/\Delta}$ Q (Major product)

(3) $\begin{array}[t]{@{}l@{}}\ \ \ \ \mathrm{CH_3\ \ CH_3}\\[-2pt]\ \ \ \ \ \ |\ \ \ \ \ \ \ |\\[-2pt]\mathrm{H_3C-C\ -\ CH-CH_3}\\[-2pt]\ \ \ \ \ \ |\\[-2pt]\ \ \ \ \mathrm{OH}\end{array}\xrightarrow{\mathrm{H^+}/\Delta}$ R (Major product)

(4) \qfigi{p2_chem_s2_q02_d}$\xrightarrow{\mathrm{H^+}/\Delta}$ S (Major product)

If the sum of $\alpha$-hydrogen in (P, Q, R, S) is equal to Z then what is the value of $\dfrac Z{10}$?
\qn{3} Molecular formula of an open chain compound (X) is C$_8$H$_{14}$. It decolourises bromine water. On reductive ozonolysis it forms acetaldehyde and butan--1, 4-dial in molar ratio 2 : 1 ?

How many allylic carbon atoms are present in one molecule of compound (X)?
\qn{4} $\mathrm{H_2C{=}CH-CH(CH_3)-CH(CH_3)_2}$.

Its R-configuration is subjected to monochlorination. How many compounds with inversion of configuration can be obtained ?
\qn{5} The total number of paramagnetic and inner orbital complex of the following is

$\mathrm{Na_4[Fe(CN)_5(NOS)]}$, $\mathrm{[Mn(NH_3)_6]\,SO_4}$, $\mathrm{[Cr(NH_3)_6]Cl_3}$

$\mathrm{K_2[PtCl_6]}$,$\mathrm{[Ni(NH_3)_6]\,SO_4}$, $\mathrm{[V(H_2O)_6]SO_4}$, $\mathrm{K_3[CuCl_4]}$
\qn{6} What should be the number for Cl in the IUPAC naming of given compound?
\qfig{p2_chem_s2_q06_1}
\qn{7} Among the following, total number of metals from which metal is extracted by self reduction method -

Pb, Zn, Cu, Mg, Hg, Fe, Al, Au
\qn{8} By passing a certain amount of charge through NaCl solution, 9.08 L of Chlorine gas were liberated at STP. When the same amount of charge is passed through a nitrate solution of a metal, 52.8 gm, of the metal was deposited. If specific heat of metal is 0.032 cal/$^\circ$C-g, the valency of metal is:
\qn{9} A certain mass of anhydrous oxalic acid is converted into H$_2$O, CO$_2$ and CO, on heating in presence of H$_2$SO$_4$. The CO formed reacts completely with iodine pentoxide to liberate iodine and CO$_2$. The iodine thus liberated required 200 ml of 0.2 N thiosulphate for end point. The mass (in g) of oxalic acid taken was:
\qn{10} An aqueous solution prepared has initially 0.5M Fe$^{+3}$(aq.) and 1M KI(aq.). If molarity of I$^-$ ions at equilibrium is $(x\times10^{-5})$ then value of $x$ is ?

(Given $E^\circ_{\mathrm{Fe^{3+}|Fe^{2+}}}=0.77$V, $E^\circ_{\mathrm{I_2|I^-}}=0.53$V, $2.303\,RT/F=0.06$)
\partsec{3}{MATHEMATICS}
\sect{Section-I}
\qn{1} Value of $\dfrac{4\sin9^\circ\,\sin21^\circ\,\sin39^\circ\,\sin51^\circ\,\sin69^\circ\,\sin81^\circ}{\sin54^\circ}$ is equal to
\begin{qoptions}
\opt{A}{$\dfrac1{16}$}
\opt{B}{$\dfrac1{32}$}
\opt{C}{$\dfrac18$}
\opt{D}{$\dfrac14$}
\end{qoptions}
\qn{2} The equation of the directrix of the parabola $y^2+4y+4x+2=0$ is -
\begin{qoptions}
\opt{A}{$x=-1$}
\opt{B}{$x=1$}
\opt{C}{$x=-\dfrac32$}
\opt{D}{None of these}
\end{qoptions}
\qn{3} The length of the normal (terminated by the major axis) at a point (not being the vertices) of the ellipse $\dfrac{x^2}{a^2}+\dfrac{y^2}{b^2}=1$ is $(a>b>0)$ -

(where $r$ and $r_1$ are the focal distance of the point.)
\begin{qoptions}
\opt{A}{$\dfrac ba(r+r_1)$}
\opt{B}{$\dfrac ba|r-r_1|$}
\opt{C}{$\dfrac ba\sqrt{rr_1}$}
\opt{D}{independent of $r$, $r_1$}
\end{qoptions}
\qn{4} Suppose that, $\log_{10}(x-2)+\log_{10}y=0$ and $\sqrt x+\sqrt{y-2}=\sqrt{x+y}$. Then the value of $(x+y)$, is
\begin{qoptions}
\opt{A}{2}
\opt{B}{$2\sqrt2$}
\opt{C}{$2+2\sqrt2$}
\opt{D}{$4+2\sqrt2$}
\end{qoptions}
\qn{5} AB is a double ordinate of the hyperbola $\dfrac{x^2}{a^2}-\dfrac{y^2}{b^2}=1$ such that $\Delta$AOB (where `O' is the origin) is an equilateral triangle, then the eccentricity $e$ of the hyperbola satisfies -
\begin{qoptions}
\opt{A}{$e>\sqrt3$}
\opt{B}{$1<e<\dfrac2{\sqrt3}$}
\opt{C}{$e<\dfrac2{\sqrt3}$}
\opt{D}{$e>\dfrac2{\sqrt3}$}
\end{qoptions}
\qn{6} Number of values of `$x$' in $(-2\pi,2\pi)$ satisfying the equation $2^{\sin^2x}+4\times2^{\cos^2x}=6$ is -
\begin{qoptions}
\opt{A}{8}
\opt{B}{6}
\opt{C}{4}
\opt{D}{2}
\end{qoptions}
\qn{7} Two tangents to the parabola $y^2=4ax$ make angles $\alpha_1$ and $\alpha_2$ with the positive direction of $x$-axis. The locus of their point of intersection if $\dfrac{\cot\alpha_1}{\cot\alpha_2}=2$ is :
\begin{qoptions}
\opt{A}{$2y^2=9ax$}
\opt{B}{$4y^2=9ax$}
\opt{C}{$y^2=9ax$}
\opt{D}{none of these}
\end{qoptions}
\qn{8} Number of integral solution of the equation, $4\log_{x/2}\left(\sqrt x\right)+2\log_{4x}\left(x^2\right)=3\log_{2x}\left(x^3\right)$ is :
\begin{qoptions}
\opt{A}{0}
\opt{B}{1}
\opt{C}{2}
\opt{D}{None of these}
\end{qoptions}
\qn{9} The line, $\ell x+my+n=0$ will cut the ellipse $\dfrac{x^2}{a^2}+\dfrac{y^2}{b^2}=1$ in points whose eccentric angles differ by $\pi/2$ if -
\begin{qoptions}
\opt{A}{$a^2\ell^2+b^2n^2=2m^2$}
\opt{B}{$a^2m^2+b^2\ell=2n^2$}
\opt{C}{$a^2\ell^2+b^2m^2=2n^2$}
\opt{D}{$a^2n^2+b^2m^2=2\ell$}
\end{qoptions}
\qn{10} The locus of the foot of the perpendicular from the centre of the hyperbola $xy=c^2$ on a variable tangent is :
\begin{qoptions}
\opt{A}{$(x^2-y^2)^2=4c^2xy$}
\opt{B}{$(x^2+y^2)^2=2c^2xy$}
\opt{C}{$(x^2-y^2)=4c^2xy$}
\opt{D}{$(x^2+y^2)^2=4c^2xy$}
\end{qoptions}
\qn{11} If A,B,C are the angles of a triangle such that $A+B\in\left(0,\dfrac\pi2\right]$ satisfying $\sin^2A+\sin^2B=\sin C$ then
\begin{qoptions}
\opt{A}{$C\in(0,30^\circ)$}
\opt{B}{$C\in(30^\circ,60^\circ)$}
\opt{C}{$C\in(60^\circ,90^\circ)$}
\opt{D}{$C=90^\circ$}
\end{qoptions}
\qn{12} If the line $y-\sqrt3x+3=0$ cut the parabola $y^2=x+2$ at A and B, then the value of $|PA|\cdot|PB|$ is [where $P=\left(\sqrt3,0\right)$]
\begin{qoptions}
\opt{A}{$\dfrac{4\left(2+\sqrt3\right)}{3}$}
\opt{B}{$\dfrac{4\left(2+\sqrt3\right)}{\sqrt3}$}
\opt{C}{$\dfrac{4\left(2-\sqrt3\right)}{3}$}
\opt{D}{$\dfrac{4\left(2-\sqrt3\right)}{\sqrt3}$}
\end{qoptions}
\qn{13} Let the double ordinate PP$'$ of the hyperbola $\dfrac{x^2}{4}-\dfrac{y^2}{3}=1$ is produced both sides to meet asymptotes of hyperbola in Q and Q$'$. The product (PQ)(PQ$'$) is equal to
\begin{qoptions}
\opt{A}{3}
\opt{B}{4}
\opt{C}{1}
\opt{D}{5}
\end{qoptions}
\qn{14} If $\log_{10}x+\log_{10}y=2$ then the smallest possible value of $(x+y)$ is
\begin{qoptions}
\opt{A}{10}
\opt{B}{30}
\opt{C}{20}
\opt{D}{None of these}
\end{qoptions}
\qn{15} The curve represented by the equation $\dfrac{x^2}{\sin\sqrt2-\cos\sqrt3}+\dfrac{y^2}{\sin\sqrt3-\cos\sqrt2}=1$, is
\begin{qoptions}
\opt{A}{an ellipse with foci on $x$-axis}
\opt{B}{an ellipse with foci on $y$-axis}
\opt{C}{hyperbola with foci on $x$-axis}
\opt{D}{hyperbola with foci on $y$-axis}
\end{qoptions}
\qn{16} Minimum distance between P$(a,b)$ and Q$(c,d)$ if $b^2-4a-8b+40=0$ and $c^2-8c-4d+40=0$, is
\begin{qoptions}
\opt{A}{0}
\opt{B}{$\sqrt3$}
\opt{C}{$2\sqrt2$}
\opt{D}{$\sqrt2$}
\end{qoptions}
\qn{17} In a $\Delta$PQR, $\cos3P+\cos3Q+\cos3R=1$ then which of the following is correct
\begin{qoptions}
\opt{A}{$\Delta$PQR is acute angled triangle}
\opt{B}{$\Delta$PQR is obtuse angled triangle}
\opt{C}{$\Delta$PQR is right angled triangle}
\opt{D}{$\Delta$PQR is equilateral}
\end{qoptions}
\qn{18} The chord PQ of the rectangular hyperbola $xy=a^2$ meets the axis of $x$ at A, C is mid point of PQ and `O' is the origin. Then the $\Delta$ACO is :
\begin{qoptions}
\opt{A}{equilateral}
\opt{B}{isosceles}
\opt{C}{right angled}
\opt{D}{right isosceles}
\end{qoptions}
\qn{19} If the angle of intersection between ellipse $x^2+4y^2=4$ and circle $x^2+y^2-2=0$ is $\theta$, then $\tan\theta$ is equal to
\begin{qoptions}
\opt{A}{$\dfrac12$}
\opt{B}{$\dfrac1{\sqrt2}$}
\opt{C}{$\dfrac1{2\sqrt2}$}
\opt{D}{$\dfrac23$}
\end{qoptions}
\qn{20} Let a ray of light be travelling along the line $4x+3y-2k=0$ and strikes a parabolic mirror at $(1,p)$. After reflection it moves along the line $3x-4y+k=0$. If the directrix of the parabola is the line $3x-4y+15=0$ then the length of latus rectum of the parabola is :
\begin{qoptions}
\opt{A}{4}
\opt{B}{6}
\opt{C}{8}
\opt{D}{16}
\end{qoptions}
\sect{Section-II}
\qn{1} If the tangent and normal at a point on rectangular hyperbola $xy=c^2$ cut-off intercept $a_1$, $a_2$ on $x$-axis and $b_1$, $b_2$ on the $y$-axis, then $a_1a_2+b_1b_2$ is equal to :
\qn{2} If $\log_2\left(\log_2(\log_3x)\right)=\log_2\left(\log_3(\log_2y)\right)=0$, then the value of $(x+y)$.
\qn{3} Let $d$ be the perpendicular distance from the centre of the ellipse $\dfrac{x^2}{a^2}+\dfrac{y^2}{b^2}=1$ to the tangent drawn at a point P on the ellipse. If F$_1$ \& F$_2$ are the two foci of the ellipse, then if that $(PF_1-PF_2)^2=ma^2\left[1-\dfrac{b^2}{d^2}\right]$. Then the value of $m$.
\qn{4} Let L$_1$ : $x+y=0$ and L$_2$ : $x-y=0$ be tangents to a parabola whose focus is S(1, 2). If the length of the latus rectum of the parabola can be expressed as $\dfrac m{\sqrt n}$ (where $m$ and $n$ are co-prime), then the value of $(m+n)$.
\qn{5} If $(1+\sin t)(1+\cos t)=\dfrac54$. Find the value of $(1-\sin t)(1-\cos t)+\sqrt{10}$.
\qn{6} The line $ax+by+c=0$ is normal to $xy=1$, then $\left|\dfrac{|a|}a-\dfrac b{|b|}\right|$ is $(ab\ne0)$
\qn{7} Number of values of $x$ satisfying the equation $\sqrt[4]{|x-3|^{x+1}}=\sqrt[3]{|x-3|^{x-2}}$, is
\qn{8} Equation of ellipse given by $25\,(3x+4y-7)^2+9(4x-3y-1)^2=5625$ then equation of latus rectum of ellipse is $4x-3y=a$ where $a>0$ then $a$ is
\qn{9} If $\displaystyle\sum_{r=1}^{n}\left(\dfrac{\tan2^{r-1}}{\cos2^r}\right)=\tan p^n-\tan q\ \ \forall\ n\in\mathbb N$. Then the value of $(p+q)$ (Where $p,q\in\mathbb N$)
\qn{10} Abscissa of two points P and Q on parabola $y^2=8x$ are roots of equation $x^2-17x+11=0$. Tangents at P and Q meet at point T. Distance of T from the focus of parabola is
\end{multicols}
\end{document}
